\section{Grenzen der Umsetzung}
Die entwickelte WordClock stellt einen funktionsfähigen Prototypen dar.  
Dennoch bestehen verschiedene technische und methodische Grenzen.  

Durch den implementierten Boot-Modus ist ein wesentlich leistungsfähigeres Netzteil nötig, was bei der Auslegung berücksichtigt werden musste und sich auf die Herstellungskosten ausgewirkt hat.  
Auch die Genauigkeit einzelner Zusatzfunktionen unterliegt Einschränkungen.  
Die Darstellung der Mondphase erfolgt beispielsweise in vereinfachter Form mit nur sechs Zuständen.  
Dadurch können geringfügige Abweichungen zur tatsächlichen Mondphase auftreten.  

Weiterhin hängt die Verfügbarkeit einzelner Funktionen von externen Diensten und Netzwerkverbindungen ab.  
Bei fehlender Internetverbindung stehen Wetterdaten und Zeitsynchronisation nur eingeschränkt zur Verfügung.  

Methodische Grenzen ergeben sich aus dem prototypischen Charakter der Arbeit.  
Die entwickelte WordClock wurde als Einzelprototyp konzipiert.  
Aspekte einer industriellen Serienfertigung wurden nicht betrachtet.  
Ebenso erfolgte keine langfristige Untersuchung hinsichtlich Alterung, Langzeitstabilität oder Energieeffizienz im mehrmonatigen Dauerbetrieb.  

Trotz dieser Einschränkungen konnte ein funktionsfähiges und erweiterbares Gesamtsystem entwickelt werden.  
Die Arbeit bildet damit eine geeignete Grundlage für zukünftige Weiterentwicklungen und Optimierungen.